\documentclass[10pt,a4paper]{amsart}
\usepackage{amsmath,amssymb,amsthm,graphicx}
\usepackage{verbatim}
\newenvironment{reference}{\comment}{\endcomment}
%\newenvironment{reference}{}{}
\newenvironment{slogan}{\comment}{\endcomment}
\newenvironment{history}{\comment}{\endcomment}

% For commutative diagrams we use Xy-pic
\usepackage[all]{xy}

% We use 2cell for 2-commutative diagrams.
\xyoption{2cell}
% Theorem environments.
%
\theoremstyle{plain}
\newtheorem{theorem}[subsection]{Theorem}
\newtheorem{proposition}[subsection]{Proposition}
\newtheorem{lemma}[subsection]{Lemma}

\theoremstyle{definition}
\newtheorem{definition}[subsection]{Definition}
\newtheorem{example}[subsection]{Example}
\newtheorem{exercise}[subsection]{Exercise}
\newtheorem{situation}[subsection]{Situation}

\theoremstyle{remark}
\newtheorem{remark}[subsection]{Remark}
\newtheorem{remarks}[subsection]{Remarks}

\numberwithin{equation}{subsection}

% Macros
%
\def\lim{\mathop{\mathrm{lim}}\nolimits}
\def\colim{\mathop{\mathrm{colim}}\nolimits}
\def\Spec{\mathop{\mathrm{Spec}}}
\def\Hom{\mathop{\mathrm{Hom}}\nolimits}
\def\Ext{\mathop{\mathrm{Ext}}\nolimits}
\def\SheafHom{\mathop{\mathcal{H}\!\mathit{om}}\nolimits}
\def\SheafExt{\mathop{\mathcal{E}\!\mathit{xt}}\nolimits}
\def\Sch{\mathit{Sch}}
\def\Mor{\mathop{\mathrm{Mor}}\nolimits}
\def\Ob{\mathop{\mathrm{Ob}}\nolimits}
\def\Sh{\mathop{\mathit{Sh}}\nolimits}
\def\NL{\mathop{N\!L}\nolimits}
\def\CH{\mathop{\mathrm{CH}}\nolimits}
\def\proetale{{pro\text{-}\acute{e}tale}}
\def\etale{{\acute{e}tale}}
\def\QCoh{\mathit{QCoh}}
\def\Ker{\mathop{\mathrm{Ker}}}
\def\Im{\mathop{\mathrm{Im}}}
\def\Coker{\mathop{\mathrm{Coker}}}
\def\Coim{\mathop{\mathrm{Coim}}}

% Boxtimes
%
\DeclareMathSymbol{\boxtimes}{\mathbin}{AMSa}{"02}

%
% Macros for moduli stacks/spaces
%
\def\QCohstack{\mathcal{QC}\!\mathit{oh}}
\def\Cohstack{\mathcal{C}\!\mathit{oh}}
\def\Spacesstack{\mathcal{S}\!\mathit{paces}}
\def\Quotfunctor{\mathrm{Quot}}
\def\Hilbfunctor{\mathrm{Hilb}}
\def\Curvesstack{\mathcal{C}\!\mathit{urves}}
\def\Polarizedstack{\mathcal{P}\!\mathit{olarized}}
\def\Complexesstack{\mathcal{C}\!\mathit{omplexes}}
% \Pic is the operator that assigns to X its picard group, usage \Pic(X)
% \Picardstack_{X/B} denotes the Picard stack of X over B
% \Picardfunctor_{X/B} denotes the Picard functor of X over B
\def\Pic{\mathop{\mathrm{Pic}}\nolimits}
\def\Picardstack{\mathcal{P}\!\mathit{ic}}
\def\Picardfunctor{\mathrm{Pic}}
\def\Deformationcategory{\mathcal{D}\!\mathit{ef}}
\usepackage{cleveref}
\crefname{originalreference}{original reference}{original references}
\Crefname{originalreference}{Original reference}{Original references}
\makeatletter
\newlabel{definition-weil-cohomology-theory-classical}{{0FGV}{source}{}{sourcepointer-1}{}}
\newlabel{definition-weil-cohomology-theory-classical@cref}{{[originalreference][0][]0FGV}{[1][0][]source}{}{}{}}
\makeatother
\crefname{originalreference}{original reference}{original references}
\Crefname{originalreference}{Original reference}{Original references}
\makeatletter
\newlabel{lemma-degrees-cycles-classical}{{0FGW}{source}{}{sourcepointer-2}{}}
\newlabel{lemma-degrees-cycles-classical@cref}{{[originalreference][0][]0FGW}{[1][0][]source}{}{}{}}
\makeatother
\title{Multiplicativity of the degree map for a classical Weil cohomology theory}
\author{The Stacks Project Authors}
\date{stacks-project:weil.tex:local-snapshot-20261012; GNU FDL 1.2 or later}
\begin{document}
\maketitle
\noindent\textit{Editorial note.} All delivered mathematical text is the author's own, copied from the pinned Stacks Project chapter file weil.tex; no statement or proof step was rewritten, reordered or removed.\par\medskip
\noindent\textit{Editorial note.} The counted proof is the author's complete argument: the identification of the one dimensional top cohomology of the product, and the comparison on the class of a product of closed points.\par\medskip
\noindent\textit{Editorial note.} The label lemma-trace-product-classical names the counted statement, and the proof environment immediately following it in the source is delivered verbatim as its proof.\par\medskip
\noindent\textit{Editorial note.} The counted statement is carried by the author's own lemma environment in the chapter source.\par\medskip
\noindent\textit{Editorial note.} Pre-open gate: python operations/agent-authoring/tools/precard\_gate.py --source weil.tex --label lemma-trace-product-classical --start 1718 --end 1737 --local-id 28162 returned RESULT=OK (L28162 free, occupancy FREE, public interval clear). Pre-assembly tag self-check: canonical tag resolved and not present in reference/public-items.jsonl.\par\medskip
\noindent\textit{Editorial note.} The statement and proof are copied from the pinned chapter file; the cited inputs are kept as external pointers with their canonical Stacks tags.\par\medskip
\noindent\textit{Allowed standard prerequisites (editorial).}\begin{itemize}
\item The top cohomology of a product is one dimensional. Axiom (B) identifying the cohomology of a product with the tensor product, and the axiom making the top cohomology of an equidimensional variety one dimensional. Shows that the target of the degree map of the product is one dimensional and generated by the class of a product of closed points.
\item The degree of the class of a closed point is one. Same chapter, the classical version of the commutativity of the degree map with the degree of zero cycles. Computes the degree of the class of a product point as one.
\item Multiplicativity of the cycle class map for exterior products. Axioms (C)(a) and (C)(c) on compatibility with pullbacks and with exterior products. Expresses the class of a product point as the exterior product of the classes of the factors.
\end{itemize}\medskip
\ifdefined\hypertarget\hypertarget{sourcepointer-1}{}\fi
\noindent\textit{Source reference.} Original reference 0FGV: sources/stacks-weil/pinned-original-source.tex, source lines 1682--1682. Prerequisite pointer invoked at proof line(s) 1721; it names a standard fact of the same chapter that the counted argument uses, and every step of the counted proof built on it is retained verbatim in the delivered body.\par\medskip
\ifdefined\hypertarget\hypertarget{sourcepointer-2}{}\fi
\noindent\textit{Source reference.} Original reference 0FGW: sources/stacks-weil/pinned-original-source.tex, source lines 1695--1695. Prerequisite pointer invoked at proof line(s) 1730; it names a standard fact of the same chapter that the counted argument uses, and every step of the counted proof built on it is retained verbatim in the delivered body.\par\medskip
\begin{lemma}
\label{lemma-trace-product-classical}
Let $H^*$ be a classical Weil cohomology theory
(Definition \ref{definition-weil-cohomology-theory-classical}).
Let $X$ and $Y$ be smooth projective varieties.
Then $\int_{X \times Y} = \int_X \otimes \int_Y$.
\end{lemma}

\begin{proof}
Say $\dim(X) = d$ and $\dim(Y) = e$. By axiom (B) we have
$H^{2d + 2e}(X \times Y) = H^{2d}(X) \otimes H^{2e}(Y)$
and by axiom (A)(d) this is $1$-dimensional.
By Lemma \ref{lemma-degrees-cycles-classical}
this $1$-dimensional vector space generated by the
class $\gamma([x \times y])$ of a closed point $(x, y)$ and
$\int_{X \times Y} \gamma([x \times y]) = 1$.
Since $\gamma([x \times y]) = \gamma([x]) \otimes \gamma([y])$
by axioms (C)(a) and (C)(c) and since $\int_X \gamma([x]) = 1$ and
$\int_Y \gamma([y]) = 1$ we conclude.
\end{proof}

\end{document}
